\chapter{لماذا كيان؟}
\label{ch:why}

تبدأ من يوم عمل تعرفه جيداً، وتنتهي وأنت تعرف ما الذي يغيّره كيان فيه، ولماذا يُجمع عمل المؤسسة كله في مكان واحد.

\section{يوم عمل عادي}
\label{sec:why-day}
\begin{rcslide}{يوم عمل عادي}
\sphase{المشكلة}
\slead{في يوم واحد يتوزع عملك على أماكن كثيرة، ولا يعرف أيٌّ منها شيئاً عن الآخر.}
\scard{paper}{الكتب الرسمية}{دفتر وارد وصادر، ولا تعرف أين وصل الكتاب إلا بالسؤال.}
\scard{chat}{المهام}{تكليفات في رسائل واتساب تضيع بين المحادثات.}
\scard{sheet}{الاستمارات}{ملفات Excel بنسخ كثيرة، لا تُجمع إلا يدوياً.}
\scard{people}{شؤون الموظفين}{برنامج مستقل لا يتصل بالعمل اليومي.}
\stake{كل أداة تؤدي عملها وحدها، لكن المعلومة تتجزأ بينها. وكلما زادت الأماكن زاد ما يضيع.}
\end{rcslide}

\section{سؤال بسيط من المدير}
\label{sec:why-question}
\begin{rcslide}{سؤال بسيط من المدير}
\sphase{المشكلة}
\slead{سؤال واحد في اجتماع الصباح، ولكي تجيب عنه تبحث في ثلاثة أماكن على الأقل.}
\squote{كم كتاباً ورد إلينا هذا الشهر ولم نردّ عليه بعد؟ ومن المسؤول عن كل كتاب؟}
\schain{دفتر الوارد ← محادثات الزملاء ← جدول المتابعة ← جدول جديد للمدير}
\stake{يستغرق الجواب يوماً أو يومين، وقد يصل ناقصاً. والمشكلة ليست في الموظف، بل في معلومات لا يربطها شيء.}
\end{rcslide}

\section{ما هو كيان؟}
\label{sec:why-what}
\begin{rcslide}{ما هو كيان؟}
\sphase{الحل}
\slead{منصة واحدة يجد فيها كل موظف عمله، ويعرف فيها المدير من المسؤول عن كل عمل وأين وصل.}
\scard{tasks}{تطبيقات جاهزة}{تبدأ بها من أول يوم.}
\scard{form}{أدوات بناء}{تصمم بها ما يخص مؤسستك دون برمجة.}
\scard{lock}{صلاحيات دقيقة}{تحدد ما يراه كل شخص وما يفعله.}
\scard{ai}{ذكاء اصطناعي}{يجيبك عن بيانات المؤسسة التي يحق لك رؤيتها.}
\stake{لا يفرض كيان طريقة عمل جاهزة: ما يتوفر جاهزاً تستخدمه، وما يخص مؤسستك وحدها تبنيه.}
\slink{main}{ابدأ مع كيان — الدليل الرئيسي}{ar/c1.html}
\end{rcslide}

\section{مكان واحد، وكل شيء مترابط}
\label{sec:why-linked}
\begin{rcslide}{مكان واحد، وكل شيء مترابط}
\sphase{الحل}
\slead{حين تفتح كتاباً وارداً في كيان، لا ترى ورقة معزولة، بل حلقة في سلسلة تعرف أولها وآخرها.}
\schain{الكتاب الوارد ← الإحالة ← الموظف المكلّف ← الرد الصادر ← الاعتماد والتوقيع}
\scard{building}{الجهة المرسلة}{ومتى قُيّد الكتاب.}
\scard{people}{منصب المكلّف}{ووحدته في الهيكل التنظيمي.}
\scard{tasks}{مهامه الأخرى}{وملفه في الموارد البشرية.}
\stake{تُسجَّل كل معلومة مرة واحدة وتُربط بما حولها، فلا تبحث عنها ولا تنسخها من مكان إلى آخر.}
\slink{correspondence}{تطبيق المراسلات خطوة بخطوة}{ar/index.html}
\end{rcslide}

\section{أدوات جاهزة تبدأ بها اليوم}
\label{sec:why-ready}
\begin{rcslide}{أدوات جاهزة تبدأ بها اليوم}
\sphase{الحل}
\slead{الأعمال التي تتشابه في أغلب المؤسسات تجدها جاهزة، فلا تنتظر تطويراً خاصاً.}
\scard{chat}{التواصل}{البريد الإلكتروني، والمحادثات، والتقويم.}
\scard{tasks}{العمل اليومي}{المهام، وطلبات الموافقة، ومكتبة الملفات.}
\scard{people}{الموارد البشرية}{الهيكل والمناصب والموظفون والدوام والإجازات والرواتب.}
\scard{mail}{تطبيقات المؤسسة}{المراسلات الرسمية، والاستراتيجية والتنفيذ، ودليل الإجراءات.}
\stake{هذه التطبيقات مترابطة من أول يوم: المهمة تعرف الموظف، والموظف يعرف منصبه، والكتاب يعرف من أُحيل إليه.}
\slink{main}{الموارد البشرية}{ar/c3.html}
\slink{main}{المهام}{ar/c4.html}
\slink{main}{المشاريع والتخطيط}{ar/c5.html}
\slink{main}{المراسلات الرسمية}{ar/c7.html}
\end{rcslide}

\section{أدوات بناء لما يخص مؤسستك}
\label{sec:why-build}
\begin{rcslide}{أدوات بناء لما يخص مؤسستك}
\sphase{الحل}
\slead{ما لا تجده جاهزاً يبنيه مسؤول مخوَّل بالاختيار والتعبئة، ويعدّله كلما تغيّر عملكم.}
\scompare{ما تستخدمه اليوم}{استمارة كشف ورقية \\ دفتر المنشآت \\ «وقّع وحوّل للمدير» \\ تقرير شهري يُجمع يدوياً}{ما يصبح في كيان}{نموذج يملؤه المفتش من جواله \\ سجل لكل منشأة يجمع كشوفها \\ إجراء داخلي ينقل الطلب وحده \\ لوحة معلومات محدّثة دائماً}
\scard{entity}{الكيانات}{سجل لما تتابعه.}
\scard{form}{النماذج}{استمارة بقواعد وتحقق.}
\scard{flow}{الإجراءات الداخلية}{اعتماد وتكليف وإشعار.}
\scard{chart}{لوحات المعلومات}{أرقام عملك أمامك.}
\stake{بأدوات البناء يطابق كيان طريقة عملكم، بدل أن تغيّروا طريقة عملكم لتناسب البرنامج.}
\slink{builder}{دليل أدوات البناء}{ar/builder.html}
\slink{inspection}{تطبيق عملي: الكشف على المنشآت}{ar/index.html}
\end{rcslide}

\section{الذكاء الاصطناعي يعمل على بيانات مترابطة}
\label{sec:why-ai}
\begin{rcslide}{الذكاء الاصطناعي يعمل على بيانات مترابطة}
\sphase{الحل}
\slead{تسأل بلغتك كما تسأل زميلاً، فتحصل على تقرير مكتوب وأرقام من بيانات المؤسسة.}
\squote{كم زيارة أُلغيت هذا الشهر، ولماذا؟}
\scard{ai}{تقارير الذكاء الاصطناعي}{سؤال وجواب للأسئلة السريعة.}
\scard{chart}{لوحات الذكاء الاصطناعي}{تصف اللوحة التي تريدها فتُبنى لك.}
\scard{lock}{ضمن صلاحياتك}{لا يقرأ إلا ما يحق لك رؤيته، ويفعّله مسؤول المؤسسة.}
\stake{لا يجيب الذكاء الاصطناعي جيداً إلا عن بيانات يراها كاملة ومترابطة. ولأن عمل المؤسسة كله في كيان، يجد الجواب في مكان واحد.}
\slink{builder}{تقارير الذكاء الاصطناعي}{ar/builder-ai-reports.html}
\slink{builder}{لوحات معلومات الذكاء الاصطناعي}{ar/builder-ai-dashboards.html}
\end{rcslide}

\section{السؤال نفسه، بعد كيان}
\label{sec:why-after}
\begin{rcslide}{السؤال نفسه، بعد كيان}
\sphase{الفرق}
\slead{تفتح الوارد، وتصفّي الكتب غير المنجزة لهذا الشهر، فترى أمام كل كتاب حالته والمكلّف به. أو تسأل الذكاء الاصطناعي السؤال كما قاله المدير.}
\smetric{مدة الجواب}{يوم أو يومان}{دقائق}
\smetric{مصدر الجواب}{ذاكرة الزملاء}{السجل نفسه}
\smetric{دقة الجواب}{بحسب من جمعه}{واحدة للجميع}
\smetric{عدد البرامج}{أربعة أو خمسة}{برنامج واحد}
\stake{هذه هي الفائدة الكبرى من كيان: عمل المؤسسة كله في مكان واحد ومترابط، فأي نتيجة أو تقرير تحتاجه تصل إليه من هنا.}
\end{rcslide}

\section{مرن لأي مؤسسة}
\label{sec:why-any}
\begin{rcslide}{مرن لأي مؤسسة}
\sphase{الفرق}
\slead{البداية واحدة، وتأخذ كل مؤسسة منها ما يناسبها، ثم تبني الباقي.}
\scard{gov}{وزارة أو مديرية}{مراسلات رسمية بتواقيع وأختام، وجولات تفتيش، وخدمات للمواطنين.}
\scard{heart}{منظمة أو جمعية}{مشاريع بأطر منطقية وميزانيات، ومستفيدون، وتقارير للمانحين.}
\scard{building}{أي مؤسسة}{موظفون ودوام ومهام وطلبات داخلية واعتمادات.}
\schain{الهيكل والمناصب ← النماذج ← الإجراءات ← التقارير}
\stake{لا يحتاج التحول الرقمي في كيان إلى مشروع برمجي طويل؛ تنقل عملك إليه خطوة خطوة، بالترتيب أعلاه.}
\slink{inspection}{مثال حكومي: الكشف على المنشآت}{ar/index.html}
\slink{correspondence}{مثال حكومي: المراسلات}{ar/index.html}
\slink{main}{مثال للمنظمات: المستفيدون في المشاريع}{ar/projects-beneficiaries.html}
\end{rcslide}

\section{ما الذي ستتعلمه في هذه الحقيبة؟}
\label{sec:why-next}
\begin{rcslide}{ما الذي ستتعلمه في هذه الحقيبة؟}
\sphase{الخطوة التالية}
\slead{كل وحدة قصيرة، وكل شريحة فيها فكرة واحدة. والتفاصيل في الأدلة المختصة، على بعد رابط.}
\scard{chat}{أدوات العمل اليومي والتواصل}{البريد والمحادثات والتقويم والملفات والمهام والموافقات.}
\scard{building}{تطبيقات المؤسسة}{الموارد البشرية والاستراتيجية والمراسلات والإجراءات والإعدادات.}
\scard{form}{أدوات البناء}{الكيانات والنماذج والإجراءات ولوحات المعلومات.}
\scard{paper}{من الورق إلى كيان}{كيف ينتقل عمل مؤسستك إليه، مع تمرين عملي.}
\scard{ai}{التقارير والذكاء الاصطناعي}{كيف تصل إلى أي نتيجة تحتاجها، ثم تقييم ختامي.}
\stake{أنت الآن تعرف لماذا تتعلم كيان. انتقل إلى الوحدة التالية: أدوات العمل اليومي والتواصل.}
\end{rcslide}
